
% ---------------------------------------------------------------------------
\section{Computer-assisted steps in the proof of Proposition~\ref{prop:Psigma-unique}}\label{app:certificates}

This appendix describes how Lemmas~\ref{lem:Psigma-convex} and~\ref{lem:Psigma-inclusion} are established and states, step by step, which parts are proved by hand and which rest on interval arithmetic. The computations are carried out with ball arithmetic (Arb~\cite{johansson2017arb}, accessed through python-flint~\cite{pythonflint}) and, for some steps, with interval arithmetic in mpmath~\cite{mpmath}. Here ``certified'' means: the parameter box is covered by finitely many sub-boxes, and on each one a rigorously computed lower bound is positive. The soundness of the libraries is assumed; nothing in this appendix is formally verified. Two analytic auxiliary estimates needed for $p\ge23$ (the weights of the guaranteed divisors and the remainder of the cutoff terms) are proved by hand in Lemmas~\ref{lem:class-weights} and~\ref{lem:remainder} below. The scripts, the logs of the runs quoted here and a short \texttt{README} are in the repository directory \texttt{src/certificates}.

\subsection{A representation by poles}
For $0\le t<1$ and $\rho\in\mathbb Z$ put
\[
g_\rho(t):=\frac{\sin^2(\pi t)}{\pi^2(\rho+t)^2},\qquad \sum_{\rho\in\mathbb Z}g_\rho(t)=1
\]
(the partial-fraction expansion of $\csc^2$ at $z=t$). By Proposition~\ref{prop:RPF}, for $i\ge2$ and $x=p+t$,
\[
\frac{F(x,i)}{i^2}=\sum_{n\in\{ik:k\in\mathbb Z\}}g_{p-n}(t).
\]
All terms are non-negative, so the double series may be rearranged, and
\begin{equation}\label{eq:pole-rep}
\mathcal P_\sigma(p+t;\kappa)=\sum_{\rho\in\mathbb Z}W(\rho,t)\,g_\rho(t)-p-t,\qquad
W(\rho,t):=\sum_{\substack{i\ge2\\ i\mid(p-\rho)}}i\,\phi_\kappa\!\Bigl(\frac{i}{p+t+1}\Bigr),
\end{equation}
where for $\rho=p$ every $i\ge2$ divides $p-\rho=0$. For $\kappa=\infty$ the weights become $w_p(\rho):=\sum_{i\mid(p-\rho),\,2\le i\le p+1}i$, and for $0<t<1$ the limit function $\mathcal P_\infty(p+t):=\sum_{2\le i\le p+1}F(p+t,i)/i-(p+t)$ equals $\sum_\rho(w_p(\rho)-p)g_\rho(t)-t$. Only $\rho=0$ and the finitely many $\rho$ with small $|\rho|$ matter in what follows.

\subsection{Inclusion for $\kappa=\infty$ (hand-checkable part and one certificate)}
\begin{lemma}[Guaranteed divisors]\label{lem:guaranteed}
Let $p$ be an odd prime, $R\le p-4$, $0<|\rho|\le R$ and $d:=p-\rho$. Then $w_p(\rho)\ge d+\tfrac d2$ for $\rho=-1$ and for odd $\rho\ge1$; $w_p(\rho)\ge d$ for even $\rho\ge2$; $w_p(\rho)\ge\tfrac{p+|\rho|}{2}$ for odd $\rho\le-3$; and $w_p(\rho)\ge0$ for even $\rho\le-2$. Moreover $w_p(0)=p$.
\end{lemma}
\begin{proof}
$d$ and, if $d$ is even, $d/2$ are divisors of $d$ in the range $[2,p+1]$ as long as $R\le p-4$ (for $\rho\le-2$ only $d/2\le p+1$ is admissible); the parity of $d$ is that of $\rho+1$ because $p$ is odd. For $\rho=0$ the only divisor of $p$ in $[2,p+1]$ is $p$.
\end{proof}
For $0<t<1$ and $R\ge2$ the tail satisfies $\sum_{|\rho|>R}g_\rho(t)\le T_R(t):=\frac{\sin^2(\pi t)}{\pi^2}\frac{2R}{R^2-t^2}$ (compare the series with $\int_R^\infty(u\pm t)^{-2}\,du$). With $c(-1)=\tfrac12$, $c(\rho)=\tfrac12$ for odd and $0$ for even $\rho\ge1$, and $c(\rho)=-\tfrac12$ for odd and $-1$ for even $\rho\le-2$, one obtains for $p\ge R+4$
\begin{equation}\label{eq:LR}
\frac{\mathcal P_\infty(p+t)}{p}\ \ge\ L_R(p,t):=\sum_{0<|\rho|\le R}c(\rho)g_\rho(t)-\frac{3}{2p}\sum_{\rho=1}^{R}\rho\,g_\rho(t)-T_R(t)-\frac tp .
\end{equation}
For $R=19$ the only $p$-dependent terms are $-\tfrac3{2p}\sum_{\rho\le19}\rho g_\rho-\tfrac tp$, which increase with $p$. Hence $L_{19}(p,t)\ge L_{19}(23,t)$ for $p\ge23$, and the \emph{certificate} is $L_{19}(23,t)>0.0015$ on $[\tfrac{9}{50},1]$ (value $0.0018456\ldots$ at $t=\tfrac{9}{50}$, attained at the left end; verified in Arb with $200$ bit by two different covers). It follows that
\begin{equation}\label{eq:Pinf-incl}
\mathcal P_\infty(p+t)\ \ge\ 0.0015\,p>0\qquad(p\ge23,\ \tfrac9{50}\le t<1).
\end{equation}
For $p\in\{5,7,11,13,17,19\}$ the exact integer weights $w_p(\rho)$ are used in a window $|\rho|\le300$ and the same tail bound outside; the resulting lower bounds on $[\tfrac9{50},1]$ are certified in the same way.

\subsection{Inclusion for finite $\kappa$ (Lemma~\ref{lem:Psigma-inclusion})}\label{subsec:finite-kappa}
Write $f_i:=F/i^2$ and split, for $0<t<1$,
\[
\mathcal P_\sigma=\mathcal P_\infty-(1-\phi_{p+1})E-D+H,\quad
E:=(p+1)f_{p+1},\ \ D:=\sum_{2\le i\le p}i(1-\phi_i)f_i,\ \ H:=\sum_{i\ge p+2}i\,\phi_if_i,
\]
with $\phi_i:=\phi_\kappa(i/(p+t+1))$ and $D\ge0$, $H\ge0$. Hence $\mathcal P_\sigma\ge\mathcal P_\infty-(1-\phi_{p+1})E-D$.

\emph{Reduction to $p=23$.} Let $p\ge23$, $\kappa\ge C(p+1)\ln p$ and $t\in[t_{\min},1)$, where $t_{\min}:=\tfrac{23}{128}$ is a dyadic number slightly below $\tfrac9{50}$ (used for the interval cover). Write $\operatorname{sinc}y:=\sin(\pi y)/(\pi y)$ and $g_n(s):=\bigl(\operatorname{sinc}s/\operatorname{sinc}(s/n)\bigr)^2$ for $0<s\le1$. Three elementary estimates give
\[
\frac{(1-\phi_{p+1})E+D}{p}\ \le\ \Delta_C(t):=\tfrac{24}{23}\,g_{24}(1-t)\,W_C(t)+D_0(C),
\]
with $W_C(t):=\exp\bigl(-2C\ln23\cdot t/(1+t/24)\bigr)$, $D_0(C):=e^{-a_0(1+t_{\min})}/(1-e^{-a_0})$ and $a_0:=2C\ln23\cdot\tfrac{24}{25}$:
\begin{enumerate}
\item[(i)] $1-\phi_{p+1}=\bigl(1+e^{2\kappa t/(p+1+t)}\bigr)^{-1}\le\exp\bigl(-2C\ln p\cdot t/(1+\tfrac t{p+1})\bigr)\le W_C(t)$, because $\ln p\ge\ln23$ and $1+\tfrac t{p+1}\le1+\tfrac t{24}$.
\item[(ii)] $E/p=\tfrac{p+1}{p}\,g_{p+1}(1-t)$, since $\sin\bigl(\pi(p+t)/(p+1)\bigr)=\sin\bigl(\pi(1-t)/(p+1)\bigr)$. For fixed $s$ the function $n\mapsto g_n(s)$ is decreasing, because $n\sin(\pi s/n)$ increases with $n$ ($\tan y>y$ on $(0,\tfrac\pi2)$); hence $E/p\le\tfrac{24}{23}\,g_{24}(1-t)$.
\item[(iii)] For $2\le i\le p$ put $j:=p+1-i\ge1$ and $a:=2\kappa/(p+1+t)$. Then $1-\phi_i=(1+e^{a(j+t)})^{-1}\le e^{-a(j+t)}$, and $f_i\le1$, $i\le p$ give $D/p\le\sum_{j\ge1}e^{-a(j+t)}=e^{-a(1+t)}/(1-e^{-a})$. This decreases in $a$, and $a\ge a_0$ because $(p+1)/(p+1+t)\ge(p+1)/(p+2)\ge\tfrac{24}{25}$ and $\ln p\ge\ln23$; with $t\ge t_{\min}$ it is at most $D_0(C)$.
\end{enumerate}
By \eqref{eq:Pinf-incl} we get $\mathcal P_\sigma(p+t)/p\ge0.0015-\Delta_C(t)$ for $\tfrac9{50}\le t<1$. The certificate is the one-variable inequality $\sup_{t\in[t_{\min},1]}\Delta_C(t)<0.0015$ (interval arithmetic in mpmath, repeated in Arb). It closes for $C=3.05$ with $\sup_t\Delta_{3.05}(t)\le1.4973\cdot10^{-3}$, a margin of $0.2\,\%$ below the target; the same computation closes for $C\ge3.044$ and not for $C\le3.043$, so $3.05$ is close to what this chain can give. Together with \eqref{eq:Pinf-incl} this proves $\mathcal P_\sigma(p+t)>0$ for $p\ge23$, $\tfrac9{50}\le t<1$ and $\kappa\ge3.05\,(p+1)\ln p$. For $p\le19$ the sum $\Lambda:=\sum_{2\le i\le p+1}i\,\phi_if_i-x\le\mathcal P_\sigma$ is increasing in $\kappa$, so a certificate at the smallest admissible $\kappa$ covers all larger ones. On a grid of mesh $0.05$ in $C$, the smallest values of $C$ for which the interval certificate closes are $3.25,\ 1.2,\ 0.65,\ 0.65,\ 0.5,\ 0.5$ for $p=5,7,11,13,17,19$; these are grid values, not sharp thresholds. For $p=5$ the margin is small: $f(\tfrac9{50})/p=2.5\cdot10^{-4}$ at $C=3.25$, and a direct evaluation puts the sign change of $f(\tfrac9{50})$ at $C\approx3.183$.

\subsection{Convexity (Lemma~\ref{lem:Psigma-convex})}
Differentiating \eqref{eq:pole-rep} twice gives $f''=A+B$ with
\[
A=\sum_\rho W(\rho,t)\,g_\rho''(t),\qquad B=\sum_{i\ge2}i\bigl(\phi_i''f_i+2\phi_i'f_i'\bigr).
\]
\emph{Sign of $g_\rho''$.} One computes $g_\rho''(t)=H_\rho(t)/(\pi^2y^4)$ with $y=\rho+t$ and $H_\rho(t)=2\pi^2\cos(2\pi t)\,y^2-4\pi\sin(2\pi t)\,y+6\sin^2(\pi t)$. For $0<t\le\tfrac9{50}$ and $\rho\le-1$ all three summands are non-negative. For $\rho\ge1$, $H_\rho$ is increasing in $y$ on $y\ge1$ because $\tan(2\pi\cdot\tfrac9{50})=2.125<\pi$, so $H_\rho\ge H_{1}$, and $H_1$ is positive on $[0,\tfrac9{50}]$ (a one-variable check; $H_1(\tfrac9{50})=8.09\cdot10^{-3}$, and $H_1$ changes sign at $t\approx0.18005$, which is why the method stops at $\tfrac9{50}$). Only $\rho=0$ contributes negatively, with $g_0''<0$ on $[0,\tfrac9{50}]$ (certified bound $g_0''\le-4.80$).

\emph{Lower bound for $A$.} Fix a residue class $r\in\{1,5\}$ and put $p_5:=23$, $p_1:=31$ (the primes $p\le19$ are treated individually) and $R:=40$. For integers $q\ge1$ and $0<|\rho|\le R$ let $d:=q-\rho$,
\[
G(q,\rho):=\bigl\{d/m:\ m\in\{1,2,3,6\},\ m\mid d,\ 2\le d/m\le q+1\bigr\}
\]
(a set of divisors of $q-\rho$ in $[2,q+1]$, the \emph{guaranteed divisors}; $G=\emptyset$ if $d\le0$), and
\[
c_a(q,\rho):=\frac1q\sum_{\substack{i\in G(q,\rho)\\ i\ne q+1}}i,\qquad
c_b(q,\rho):=\begin{cases}\frac{q+1}{q}&\text{if }q+1\in G(q,\rho),\\[2pt]0&\text{otherwise.}\end{cases}
\]

\begin{lemma}[Class weights]\label{lem:class-weights}
Let $r\in\{1,5\}$, $q\equiv r\pmod 6$, $q>R+12$ and $0<|\rho|\le R$. Then
\[
c_a(q,\rho)=\Bigl(1-\frac\rho q\Bigr)K_r(\rho),\qquad K_r(\rho):=\sum_{\substack{m\mid\gcd(r-\rho,\,6)\\ m\ge2\ \text{or}\ \rho\ge1}}\frac1m ,
\]
and $c_b(q,-1)=1+\tfrac1q$, $c_b(q,\rho)=0$ for $\rho\ne-1$. Consequently $q\mapsto c_a(q,\rho)$ is increasing for $\rho>0$ and decreasing for $\rho<0$, with limit $K_r(\rho)$. The exact rational numbers
\[
c_a^{(r)}(\rho):=\min\Bigl(\min_{\substack{q\equiv r\ (6)\\ p_r\le q\le300}}c_a(q,\rho),\ \ K_r(\rho)\ \text{ if }\rho<0\Bigr),\qquad
c_b^{(r)}(\rho):=\begin{cases}1&\rho=-1,\\0&\text{otherwise,}\end{cases}
\]
satisfy, for every prime $p\equiv r\pmod 6$ with $p\ge p_r$ and every $0<|\rho|\le R$,
\[
\sum_{\substack{i\in G(p,\rho)\\ i\ne p+1}}i\ \ge\ p\,c_a^{(r)}(\rho),\qquad c_b(p,\rho)\ \ge\ c_b^{(r)}(\rho).
\]
\end{lemma}
\begin{proof}
As $q\equiv r\pmod6$, $\gcd(q-\rho,6)=\gcd(r-\rho,6)$ depends only on $r$ and $\rho$; so the set of admissible $m$ (the divisors of this gcd) is the same for all such $q$. For such $m$ we have $d/m\ge(q-R)/6>2$ because $q>R+12$, so the lower bound $d/m\ge2$ is automatic. The upper bound $d/m\le q+1$ holds automatically for $m\ge2$, since $d/m\le(q+R)/2<q+1$, and for $m=1$ it holds if and only if $\rho\ge-1$. Moreover $d/m=q+1$ happens only for $m=1$, $\rho=-1$. Hence $G(q,\rho)=\{(q-\rho)/m:\ m\mid\gcd(r-\rho,6),\ m\ge2\text{ or }\rho\ge-1\}$, and $q+1\in G(q,\rho)$ exactly when $\rho=-1$. Summing $(q-\rho)/m$ over the $m$ that remain once $(m,\rho)=(1,-1)$ is removed and dividing by $q$ gives the formula for $c_a$; the formula for $c_b$ is the case $(m,\rho)=(1,-1)$. The factor $1-\rho/q$ increases in $q$ for $\rho>0$ and decreases for $\rho<0$, with limit $1$.

Now let $p\equiv r\pmod 6$ be a prime, $p\ge p_r$. If $p\le300$, then $p$ is one of the $q$ in the minimum defining $c_a^{(r)}(\rho)$. If $p>300$, let $q'$ be the largest integer $\equiv r\pmod6$ not exceeding $300$; then $q'>R+12$, and for $\rho>0$ we get $c_a(p,\rho)\ge c_a(q',\rho)\ge c_a^{(r)}(\rho)$ by monotonicity, while for $\rho<0$ we get $c_a(p,\rho)\ge K_r(\rho)\ge c_a^{(r)}(\rho)$. Finally, $c_b(p,-1)=(p+1)/p\ge1$ because $p+1\in G(p,-1)$ ($m=1$, $d=p+1$), and $c_b^{(r)}(\rho)=0$ for $\rho\ne-1$ is trivial.
\end{proof}

The finite minimum is a computation with exact rational numbers (about $2\cdot80$ numbers $c_a^{(r)}(\rho)$; the table is in \texttt{src/certificates}). The formulas of Lemma~\ref{lem:class-weights} were also compared with the definition of $c_a,c_b$ for all $q\equiv r\pmod 6$, $53\le q\le2\cdot10^5$, and the bounds with all primes $p_r\le p\le6\cdot10^6$.

\begin{lemma}[Lower bound for $A$]\label{lem:A-bound}
Let $r\in\{1,5\}$, $p\ge p_r$ a prime with $p\equiv r\pmod6$, $\kappa\ge3.25\,(p+1)\ln p$ and $0\le t\le\tfrac9{50}$. With $\sigma_0:=\sigma(6.5\ln p_r)$ and $s:=2\kappa t/(p+1+t)$,
\[
\frac Ap\ \ge\ A_0(t)+\sigma(s)A_1(t),\qquad
A_0:=g_0''+\sigma_0\!\!\sum_{0<|\rho|\le R}\!c_a^{(r)}(\rho)\,g_\rho'',\qquad
A_1:=\sum_{0<|\rho|\le R}\!c_b^{(r)}(\rho)\,g_\rho'' .
\]
\end{lemma}
\begin{proof}
Recall $A=\sum_\rho W(\rho,t)g_\rho''(t)$ with $W(\rho,t)=\sum_{i\ge2,\ i\mid(p-\rho)}i\,\phi_i$. All $g_\rho''$ with $\rho\ne0$ are positive on $(0,\tfrac9{50}]$ (see above) and $W\ge0$, so the terms with $|\rho|>R$ can be dropped. For $\rho=0$ the only divisor $i\ge2$ of $p$ is $p$, so $0<W(0,t)=p\phi_p\le p$; since $g_0''<0$ this term is at least $p\,g_0''$. For $0<|\rho|\le R$ the elements of $G(p,\rho)$ are distinct divisors of $p-\rho$ in $[2,p+1]$, so $W(\rho,t)\ge\sum_{i\in G(p,\rho)}i\,\phi_i$. For $i\le p$ we have $s_i:=2\kappa(1-\tfrac i{p+t+1})\ge2\kappa\tfrac{t+1}{p+t+1}\ge\tfrac{2\kappa}{p+1}\ge6.5\ln p_r$, hence $\phi_i\ge\sigma_0$, and $\phi_{p+1}=\sigma(s)$ exactly. Therefore $W(\rho,t)\ge\sigma_0\,p\,c_a(p,\rho)+\sigma(s)\,p\,c_b(p,\rho)$, and Lemma~\ref{lem:class-weights} gives the claim after division by $p$. The case $t=0$ follows by continuity.
\end{proof}

\emph{The edge term.} Put $B_{\rm edge}:=(p+1)\bigl(\phi_{p+1}''f_{p+1}+2\phi_{p+1}'f_{p+1}'\bigr)$ and $B_{\rm rest}:=B-B_{\rm edge}$. Since $\sin(\pi(p+t)/(p+1))=\sin\theta$ with $\theta:=\pi(1-t)/(p+1)$, one has $f_{p+1}(p+t)=t^2q(t)$ with $q:=\varrho^2$, $\varrho:=\pi\operatorname{Sc}(t)/((p+1)\sin\theta)$, $\operatorname{Sc}(t):=\sin(\pi t)/(\pi t)$. Moreover $\phi_{p+1}=\sigma(s)$, and with $a:=(p+1)/(p+1+t)$ one has $s'=as/t$ and $s''=-2as/(t(p+1+t))$. Inserting $f_{p+1}'=(2/t+q'/q)f_{p+1}$ and $q'/q=2\bigl(\operatorname{Sc}'/\operatorname{Sc}+\tfrac{\pi}{p+1}\cot\theta\bigr)$ gives
\begin{gather*}
\frac{B_{\rm edge}}{p+1}=U\,s^2\sigma''(s)+V\,s\,\sigma'(s),\\
U=a^2q,\qquad V=a\,q\,Y,\qquad T:=\pi t\cot(\pi t)-1,\\
Y:=4+4T+\frac{4\pi t}{p+1}\cot\theta-\frac{2t}{p+1+t}.
\end{gather*}
(The formulas were checked against direct differentiation in $200$ random cases.) Thus $B_{\rm edge}$ depends on $\kappa$ only through $s$, and $U,V$ depend on $p$ only through $u:=1/(p+1)$ and on $t$. This reduces the $\kappa$-dependence to the plane $(t,s)$, which is covered by boxes; the constraint $\kappa\ge3.25(p+1)\ln p$ enters as $s\ge2\cdot3.25\ln p\,(p+1)t/(p+1+t)$, a bound that increases with $p$, so that it suffices to use $p=p_r$.

\emph{The remainder.} $B_{\rm rest}=\sum_{i\ge2,\,i\ne p+1}i\,(\phi_i''f_i+2\phi_i'f_i')$. For $p\ge23$ it is controlled by the following lemma.

\begin{lemma}[Remainder of the cutoff terms]\label{lem:remainder}
Let $C=3.25$, $p\ge23$ a prime, $\kappa\ge C(p+1)\ln p$ and $0\le t\le\tfrac9{50}$. Put $\ell:=2C\ln p$, $z:=\ell\,\tfrac{p+2}{p+1}$, $\lambda:=\ell\,\tfrac{41}{50}\,\tfrac{p+1}{p+1+9/50}$ and
\[
\varepsilon(p):=\frac{e^{-\ell}}{1-e^{-\ell}}\Bigl(\ell^2+4\sqrt2\,\pi\ell+\frac{2\ell}{p+1}\Bigr)+\frac{2(p+2)}{p}\,e^{-\lambda}\Bigl(z^2+4\sqrt2\,\pi z+\frac{2\ell(p+2)}{(p+1)^2}\Bigr).
\]
Then $|B_{\rm rest}(t)|\le p\,\varepsilon(p)$, and $\varepsilon$ is decreasing on $[23,\infty)$. In particular $|B_{\rm rest}|/p\le\varepsilon(23)<1.14\cdot10^{-4}$ for all $p\ge23$, and $|B_{\rm rest}|/p\le\varepsilon(31)<2.52\cdot10^{-5}$ for $p\ge31$.
\end{lemma}
\begin{proof}
\emph{One term.} Let $s_i(t):=2\kappa\bigl(1-i/(p+t+1)\bigr)$, so that $\phi_i=\sigma(s_i)$, $s_i'=2\kappa i/(p+t+1)^2$, $s_i''=-4\kappa i/(p+t+1)^3$, $\phi_i'=\sigma'(s_i)s_i'$, $\phi_i''=\sigma''(s_i)s_i'^2+\sigma'(s_i)s_i''$. Since $\sigma'(s)=e^{-|s|}/(1+e^{-|s|})^2$ and $\sigma''=-\sigma'\tanh(s/2)$, we have $|\sigma'|,|\sigma''|\le e^{-|s|}$. Next, $f_i=h_i^2$ with $h_i(x)=\sin(\pi x)/(i\sin(\pi x/i))$, $x=p+t$ (note $\sin^2\pi(p+t)=\sin^2\pi t$). In the variable $y=\pi x/i$, $h_i=\sin(iy)/(i\sin y)=\tfrac1i\sum_{k=0}^{i-1}\cos\bigl((i-1-2k)y\bigr)$ is a real trigonometric polynomial of degree $i-1$ with $|h_i|\le1$. Bernstein's inequality gives $|dh_i/dy|\le i-1$, hence $|dh_i/dx|\le\pi(i-1)/i<\pi$ and $|f_i'|=2|h_i||h_i'|\le2\pi\sqrt{f_i}\le2\sqrt2\,\pi\sqrt{f_i}$ (the factor $\sqrt2$ is not needed; the certificates were computed with it). With $f_i\le1$, and with $z_i:=2\kappa i/(p+1)^2\ge s_i'$, $z_i':=4\kappa i/(p+1)^3\ge|s_i''|$ (as $t\ge0$), it follows that
\[
\bigl|i(\phi_i''f_i+2\phi_i'f_i')\bigr|\ \le\ T_i(\kappa):=i\,e^{-2\kappa m_i}\bigl(z_i^2+z_i'+4\sqrt2\,\pi z_i\bigr),
\]
where $m_i:=1-\tfrac i{p+1}$ for $i\le p$ and $m_i:=\tfrac i{p+1.18}-1$ for $i\ge p+2$; indeed $|s_i|\ge2\kappa m_i$ for $0\le t\le0.18$ (for $i\le p$: $1-\tfrac i{p+t+1}\ge1-\tfrac i{p+1}$; for $i\ge p+2$: $\tfrac i{p+t+1}-1\ge\tfrac i{p+1.18}-1$).

\emph{Monotonicity in $\kappa$.} $T_i(\kappa)$ is a sum of terms $c\,\kappa^ae^{-2m_i\kappa}$ with $a\in\{1,2\}$, which decrease for $\kappa\ge1/m_i$. Now $1/m_i\le p+1$ for $i\le p$ and $1/m_i\le(p+1.18)/0.82$ for $i\ge p+2$, and both are smaller than $\kappa_{\min}:=C(p+1)\ln p\ge10(p+1)$. Hence $T_i(\kappa)\le T_i(\kappa_{\min})$ for all $\kappa\ge\kappa_{\min}$.

\emph{Summation.} At $\kappa_{\min}$ we have $z_i=\ell i/(p+1)$ and $z_i'=2\ell i/(p+1)^2$. For $i\le p$ write $i=p+1-j$ with $j\ge1$; then $2\kappa_{\min}m_i=\ell j$, $z_i\le\ell$, $i\le p$, so $T_i(\kappa_{\min})\le p\,e^{-\ell j}\bigl(\ell^2+4\sqrt2\pi\ell+2\ell/(p+1)\bigr)$, and $\sum_{j\ge1}e^{-\ell j}=e^{-\ell}/(1-e^{-\ell})$. For $i\ge p+2$ we have $T_i(\kappa_{\min})=i\,e^{-2\kappa_{\min}m_i}(\alpha i^2+\beta i)$ with $\alpha,\beta>0$ independent of $i$, and $m_{i+1}-m_i=1/(p+1.18)$, so
\[
\frac{T_{i+1}(\kappa_{\min})}{T_i(\kappa_{\min})}\le\Bigl(\frac{i+1}i\Bigr)^3e^{-2\kappa_{\min}/(p+1.18)}\le\Bigl(\frac{26}{25}\Bigr)^3e^{-\ell\cdot24/24.18}<\frac12,
\]
because $i\ge25$, $\ell\ge\ell(23)=20.38\ldots$ (the left-hand side is below $2\cdot10^{-9}$). Hence $\sum_{i\ge p+2}T_i(\kappa_{\min})\le2\,T_{p+2}(\kappa_{\min})$, and $2\kappa_{\min}m_{p+2}=\lambda$, $z_{p+2}=z$, $z_{p+2}'=2\ell(p+2)/(p+1)^2$. Adding the two parts gives $|B_{\rm rest}|\le p\,\varepsilon(p)$.

\emph{Decrease in $p$.} Treat $p\ge23$ as a real variable, $\ell=2C\ln p$. In the first summand, $e^{-\ell}=p^{-2C}$; the factor $(1-p^{-2C})^{-1}$ decreases, $p^{-2C}\ell^k$ ($k=1,2$) decreases because $\frac{d}{dp}\ln(p^{-2C}\ell^k)=\bigl(k/\ln p-2C\bigr)/p<0$, and $\ell/(p+1)$ decreases for $p\ge5$. In the second summand $2(1+2/p)$, $z^2/\ell^2=\bigl(\tfrac{p+2}{p+1}\bigr)^2$, $z/\ell$ and $(p+2)/(p+1)^2$ decrease, and $e^{-\lambda}\ell^k$ ($k=1,2$) decreases because $\lambda=\ell w$ with $w=0.82\,(p+1)/(p+1.18)$ increasing, so $\lambda'\ge\ell'w\ge(2C/p)\cdot0.82\cdot\tfrac{24}{24.18}>5.2/p>(2/\ln23)/p\ge\frac d{dp}\ln\ell^k$. A product of positive decreasing functions decreases, so $\varepsilon$ decreases. Finally, interval arithmetic (Arb) gives $\varepsilon(23)=1.13874\ldots\cdot10^{-4}$ and $\varepsilon(31)=2.5126\ldots\cdot10^{-5}$.
\end{proof}

The bound is not sharp: evaluated directly (mpmath), $|B_{\rm rest}|$ is smaller than $p\,\varepsilon(p)$ by a factor of at least $2\cdot10^2$ in $4864$ cases with $23\le p\le10007$, $\kappa/\kappa_{\min}\in\{1,1.2,2,5\}$. For the finitely many $p\le19$ the same term bound is summed for each $p$ separately (with the maximum of $f_i$ over $0\le t\le\tfrac9{50}$ in place of $1$ for $i\ge p+2$, and the terms $i\ge p+62$ bounded by twice the first of them); this gives \emph{absolute} bounds $|B_{\rm rest}|\le\varepsilon_p$ with $\varepsilon_5=7.5\cdot10^{-2}$, $\varepsilon_7=1.7\cdot10^{-2}$, and smaller values for larger $p$. (Thus $\varepsilon_p$ for $p\le19$ is an absolute bound, whereas $\varepsilon(p)$ for $p\ge23$ is a bound relative to $p$.)

\emph{Certificate.} For $p\in\{5,\dots,19\}$ the lower bounds are certified in Arb (about $65$ boxes in $t$ times $2961$ boxes in $s$ up to $s=40$, with a separate estimate for $s\ge40$). For $p\ge23$ the weights are replaced by the exact class numbers $c_a^{(r)},c_b^{(r)}$ of Lemma~\ref{lem:class-weights} (classes $r=1$ and $r=5$), $\varepsilon(p)$ by $\varepsilon(p_r)$ (Lemma~\ref{lem:remainder}), the dependence on $p$ enters through $u=1/(p+1)\in[0,\tfrac1{p_r+1}]$ treated as a ball, and one certificate per class covers all $p\ge p_r$. All $p$-dependent ingredients are used in the monotone direction: $\sigma_0\ge\sigma(6.5\ln p_r)$, the lower bound for $s$ increases with $p$, $(p+1)/p\le(p_r+1)/p_r$ (this factor multiplies $U s^2\sigma''\le0$, where it is bounded above by $(p_r+1)/p_r$, and $V s\sigma'\ge0$, where it is simply dropped, being $\ge1$ with $V>0$ certified), and $\varepsilon(p)\le\varepsilon(p_r)$. The certified lower bounds for $f''/p$ are those listed in Lemma~\ref{lem:Psigma-convex}; they are weaker than the true values (for instance, at $\kappa=3.25\,(p+1)\ln p$ and $t=0.1$ one finds $f''/p=3.29$ for $p=5$ and $5.24$ for $p=23$, by direct differentiation of \eqref{eq:Psigma}).

\subsection{Status of the steps}
Table~\ref{tab:cert-status} lists, for each step of the proof of Proposition~\ref{prop:Psigma-unique}, how it is established.
\begin{table}[!h]
\centering
\small
\begin{tabular}{p{0.37\textwidth}p{0.55\textwidth}}
\hline
Step & Status\\
\hline
Pole representation \eqref{eq:pole-rep}, $\sum g_\rho=1$ & classical (partial fractions, Tonelli); proved\\
Guaranteed divisors, tail bound $T_R$ & elementary; proved. Weights additionally checked against exact divisor sums for all primes $23\le p\le2\cdot10^6$ and $38$ values of $\rho$\\
Certificate $L_{19}(23,t)>0.0015$ & interval arithmetic (Arb, $200$ bit); two independent covers; soundness assumed\\
Finite $\kappa$, reduction to $p=23$ & elementary monotonicity, proved in Section~\ref{subsec:finite-kappa} (steps (i)--(iii)); the resulting one-variable inequality $\sup_t\Delta_C(t)<0.0015$ by interval arithmetic (mpmath, repeated in Arb); $p\le19$ by interval arithmetic\\
Sign of $g_\rho''$, $\rho\ne0$ & elementary, plus a one-variable interval check\\
Edge term, formulas for $U,V$ & direct differentiation; proved (checked numerically); $V>0$ by interval arithmetic\\
Class weights (Lemma~\ref{lem:class-weights}), lower bound for $A$ (Lemma~\ref{lem:A-bound}) & elementary; proved. Finite part ($q\le300$) in exact rational arithmetic; additionally checked against all primes $p_r\le p\le6\cdot10^6$ and $0<|\rho|\le40$\\
Remainder bound $\varepsilon(p)$ (Lemma~\ref{lem:remainder}) & elementary (Bernstein's inequality, geometric series, monotonicity in $\kappa$ and in $p$); proved; $\varepsilon(23)$ and $\varepsilon(31)$ evaluated in Arb\\
Convexity, $p\le19$ & interval arithmetic (Arb), explicit absolute bounds $\varepsilon_p$; soundness assumed\\
Convexity, $p\ge23$ & interval arithmetic (Arb) uniform in $p$ per residue class $1,5\pmod6$, built on Lemmas~\ref{lem:class-weights}--\ref{lem:remainder}; soundness assumed\\
Assembly (Proposition~\ref{prop:Psigma-unique}) & elementary; proved\\
\hline
\end{tabular}
\caption{Status of the steps behind Proposition~\ref{prop:Psigma-unique}. The convexity certificate was reproduced by a second, independently written implementation (also AI-assisted); the inclusion for $\kappa=\infty$ by three computations (mpmath intervals, exact rational arithmetic, Arb). No step has been formally verified, and no human referee has examined either the certificates or the hand proofs of Lemmas~\ref{lem:guaranteed}--\ref{lem:remainder} and of the reduction in Section~\ref{subsec:finite-kappa}; the hand proofs were cross-read only by AI systems.}
\label{tab:cert-status}
\end{table}

\paragraph*{Reproduction.} From the directory \texttt{src/certificates/k2\_inclusion} run \texttt{E7927\_\ldots} (about one second); from \texttt{k1\_convexity} run the certificate scripts \texttt{E7828\_\ldots} ($p\le19$) and \texttt{E7833\_\ldots} ($p\ge23$) preceded by their gates \texttt{E7827\_\ldots} and \texttt{E7832\_\ldots}; from \texttt{k2prime\_inclusion} run \texttt{E7792\_\ldots} ($p\le19$) and \texttt{E7791\_\ldots} (reduction for $p\ge23$). The checks accompanying Lemmas~\ref{lem:class-weights}--\ref{lem:remainder}, the reduction steps of C.3, the edge-term formulas and the numerical statements of Remark~\ref{rem:right-zero-scope} and Section~\ref{subsec:Psigma-right} are in the directory \texttt{lemmas\_C}: a gate with deliberately wrong variants (\texttt{E8039\_\ldots}), the class weights (\texttt{E8040\_\ldots}), the convexity certificate with the exact class numbers (\texttt{E8041\_\ldots}), the remainder bound (\texttt{E8042\_\ldots}) and the remaining steps (\texttt{E8043\_\ldots}). The requirements are Python~3, python-flint~$0.9$, mpmath~$1.3$ and numpy. The file names carry experiment numbers of the author's working environment; the logs of the runs quoted in this appendix are stored beside the scripts.

\paragraph*{Acknowledgements}
Valuable discussions and helpful feedback from colleagues are gratefully acknowledged. The computations and analyses in Section~\ref{subsec:Psigma-right} and in this appendix were prepared with the assistance of an AI system (Claude, Anthropic); the content was checked by the author, who is responsible for it.
